\scp{Marsupials within biogeographical Wallacea}{14-01}
Looking at marsupials from west to east,
two species of cuscus are found widely on Sulawesi,
the \it{Ailurops ursinus} ‘bear cuscus’,
and the \it{Strigocuscus celebensis} ‘Sulawesi dwarf cuscus’.
A distinct species of \it{Ailurops melanotis} ‘Talaud bear cuscus’
is found in the Talaud Islands to the northeast of Sulawesi.
Another species, the \it{Strigocuscus pelengensis} ‘Banggai (dwarf)
cuscus’ is said to be found on Banggai
and on the nearby islands of the Sula archipelago in Maluku.
Other species such as \it{Spilocuscus maculatus} ‘spotted cuscus’,
\it{Phalanger ornatus} ‘ornate cuscus’,
\it{Phalanger orientalis} ‘(northern) common grey cuscus’,
and possibly the \it{Phalanger mimicus} ‘southern grey cuscus’,
are found across parts of Linguistic Wallacea.
The \it{Phalanger alexandrae} ‘Gebe cuscus’ is found on the island
of Gebe between Halmahera
and the Raja Ampat islands.
The \it{Phalanger gymnotis} ‘ground cuscus’ is found in the Raja
Ampat islands, the Bird’s Head and Bird’s Neck of Indonesian Papua, and in Aru.
The \it{Spilocuscus papuensis} ‘Waigeo spotted cuscus’ is found
on Waigeo Island in the Raja Ampat Islands.
All are marsupials of the family \it{Phalangeridae}.

It appears that the ‘Sulawesi dwarf cuscus’ (\it{Strigocuscus
celebensis}) and the larger ‘bear cuscus’ (\it{Ailurops ursinus}) on Sulawesi
are native to that region,
and did not get there through human transfer.
But in a development with many parallels to the discussion of
important food crops and words associated with them moving westward out of the New
Guinea region (see \srf{02-03-01-07}),
research summarised in \citet[260]{Schapper2011} notes that a
species of the marsupial cuscus (\it{Phalanger orientalis}) was
introduced by humans westward out from New Guinea into parts
of Maluku and Timor.
\citet[130--132]{Flannery1990} also notes that the spotted cuscus (\it{Spilocuscus
maculatus}) was transported as far west as Selayar Island off of South Sulawesi.
Like the food crops,
this westward dispersal by human transfer of some marsupial species
happened several thousand years before the arrival of the Austronesian
speaking peoples into the region.

Bandicoots (\it{Peramelidae spp.}) are much more limited in
their geographic distribution than cuscus,
hardly extending westward very far beyond Sahulland
and a few nearby islands.
That is, bandicoots are rarely found very far west of the Lydekker
Line, so Aru is included within their natural range (see \frf{02-01}).
The \it{Echymipera rufescens} ‘long-nosed spiny bandicoot, a.k.a.
long-nosed echymipera’ is found on the Aru Islands,
and also has been reported in the Kei Islands.
Words glossed ‘bandicoot’ have been identified in nine languages
of Aru, and two languages of Kei (see \srf{14-05}).
An elusive species, the \it{Rhynchomeles prattorum} ‘Seram bandicoot, a.k.a.
Seram long-nosed bandicoot’ has been identified in the mountains
of central Seram, only as far west as the region of the \tsc{Patakai-Manusela}
node of the \tsc{Ambon-Seram} subgroup (see \crf{ch:AS}).
There is as yet no evidence of the spread of bandicoots through human transfer.

Apart from bandicoots and cuscuses, additional kinds of marsupials
are found within Wallacea,
none of which have featured in the debates pertaining to marsupial
terms in Austronesian languages.
Several species of ‘tree kangaroo’ (\it{Dendrolagus spp}.,
sometimes referred to as ‘tree wallaby’) are found in Sahulland
throughout Indonesian Papua, with some on a few nearby islands to the west,
with a more restricted range in the target region than the bandicoot.
The \it{Dendrolagus inustus} ‘grizzled tree-kangaroo’,
the \it{Dendrolagus ursinus} ‘Ursine tree-kangaroo’,
and possibly the \it{Dendrolagus mayri} ‘Wondiwoi tree-kangaroo’,
are of particular interest,
being found in the Bird’s Head
and Bird’s Neck regions of Papua.
It may be worth noting that to the untrained eye,
tree kangaroos in Papua do not look significantly different from
the Sulawesi bear cuscus.\footnote{
		Grimes has seen several tree
		kangaroos in Indonesian Papua.}

Apart from tree kangaroos,
a few additional macropods occur within Wallacea, but outside New Guinea proper.
One species, \it{Thylogale brunii} ‘Dusky Pademelon’ occurs on
the Aru Islands and Kei Islands.
At least one species,
\it{Dorcposis muelleri} ‘Brown Dorcposis’,
occurs in the Raja Ampat Islands
\citep[60f, 64f]{LaveryFlannery2023}.

\it{Dasyuridae} (quolls) are another family of marsupials that occur within Wallacea.
Three extant species are attested on the Aru Islands: \it{Murexis
longicaudata} ‘Short-furred Dasyure’, \it{Myoictis wallacei} ‘Wallace’s Three-striped Dasyure’,
and \it{Sminthopis virginiae} ‘Red-cheeked Dunnart’.\footnote{
		An
		additional species of quoll,
		\it{Dasyurus albopunctatus} ‘New Guinea Quoll’,
		is common on New Guinea
		and known from the fossil record on the Aru Islands.}
An additional
species \it{Myoictis melas} ‘Müller’s Three-striped Dasyure’
occurs on Waigeo, Batanta and Salawati. \citep[42--48]{LaveryFlannery2023}.
\it{Dasyuridae} are the smallest marsupials present in Wallacea,
and despite their presence have received almost no attention
from linguists working in the region.
It is possible that several terms with vague glosses such as ‘k.o.
small marsupial’ may refer to \it{Dasyuridae} rather than bandicoots.

Finally, species of gliding possums also occur in Wallacea.
One species, \it{Distoechurus pennatus} ‘Feather-tailed Possum’
is known to occur on Waigeo Island
and “[{\ldots}] may be present on other land bridge islands associated
with the Vogelkop [Bird’s Head] Peninsula” \citep[56]{LaveryFlannery2023}.
Another species \it{Petaurus breviceps} ‘sugar glider’ occurs
in the Kei Islands,
Aru, Misool, and Halmahera \citep{Salas2016}.

Scattered reports from hunters indicate that cuscus numbers in
Wallacea are noticeably dwindling as their habitat shrinks due
to deforestation from logging,
farming, road-building, and mining (pers. comm. O. Seleky, L.
Hukunala, M. Gebhain, F. Gewagit, I. Lesnussa).
Logging poachers in protected forest areas further exacerbate the problem.
Only localised efforts at wildlife management,
setting areas of the jungle off limits to hunting for several
years, allow cuscus populations to recover slightly.
The trend is discouraging.